\section*{Step 258: Lindelof and the Subconvexity Extension}

\paragraph{Lindelof carrier.}
For zeta define
\[
\mu(1/2)=\limsup_{T\to\infty}
\frac{\log|\zeta(1/2+iT)|}{\log T}.
\]
The Lindelof hypothesis is \(\mu(1/2)=0\), equivalently
\[
\zeta(1/2+iT)=O_\epsilon(T^\epsilon)
\]
for every \(\epsilon>0\).  RH implies Lindelof; the converse is open.

\paragraph{Subconvexity family.}
For an automorphic \(L(s,\pi)\) with analytic conductor \(C(\pi)\), convexity
gives
\[
L(1/2,\pi)\ll C(\pi)^{1/4+\epsilon}.
\]
Subconvexity breaks the exponent \(1/4\).  Lindelof is the limiting exponent
\(0\).

\paragraph{Classification.}
This is neither a zero-location RH analogue nor a coefficient/zero
cross-correlation residual.  It is a critical-line growth-rate residual, so it
is parallel to both SCDG and the Selberg-Class Cross-Correlation Extension.

\paragraph{Candidate finding.}
\[
\textbf{Selberg-Class Subconvexity Extension.}
\]
The Selberg class admits a typed growth-rate family: convexity baseline,
subconvexity partial closure, and Lindelof exponent-zero target.

\paragraph{Evidence.}
\[
\begin{array}{c|c|c}
\text{instance} & \text{baseline} & \text{known result}\\
\hline
\zeta(1/2+it) & 1/4 & 13/84+\epsilon\ \text{(Bourgain)}\\
L(1/2,\chi) & 1/4 & 3/16+\epsilon\ \text{(Burgess)}\\
GL(2) & 1/4 & \text{subconvexity in aspects (Michel--Venkatesh)}
\end{array}
\]

\paragraph{Status.}
Candidate, verified on three subconvexity instances, corpus-pending.  No proof
of Lindelof, GRH, or RH is claimed.
